\subsection{DEFINITION OF THE SYMMETRIC ALGEBRA OF A MODULE}

DEFINITION 1. Let A be a commutative ring and M an A-module. The symmetric algebra of M, denoted by S(M), or Sym(M), or \(S_{A}(M)\), is the quotient algebra over A of the tensor algebra T(M) by the two-sided ideal \(\mathfrak{J}'\) (also denoted by \(\mathfrak{J}_{\mathbf{M}}'\) ) generated by the elements \(xy - yx = x \otimes y - y \otimes x\) of T(M), where \(x\) and \(y\) run through M.

Since the ideal \(\mathfrak{J}'\) is generated by homogeneous elements of degree 2, it is a graded ideal (II, §11, no. 3, Proposition 2); we write \(\mathfrak{J}_n' = \mathfrak{J}' \cap T^n(M)\) ; the algebra \(S(M)\) is then graded by the graduation (called canonical) consisting of the \(S^n(M) = T^n(M)/\mathfrak{J}_n'\) . Now \(\mathfrak{J}_0' = \mathfrak{J}_1' = \{0\}\) and hence \(S^0(M)\) is canonically identified with A and \(S^1(M)\) with \(T^1(M) = M\) ; we shall always make these identifications in what follows and denote by \(\phi'\) or \(\phi_M'\) the canonical injection \(M \to S(M)\) .

PROPOSITION 1. The algebra \(\mathbf{S}(\mathbf{M})\) is commutative.

By definition \(\phi'(x)\phi'(y) = \phi'(y)\phi'(x)\) for \(x, y\) in \(\mathbf{M}\) and, as the elements \(\phi'(x)\) , where \(x\) runs through \(\mathbf{M}\) , generate \(\mathbf{S}(\mathbf{M})\) , the conclusion follows from §1, no. 7.

PROPOSITION 2. Let E be an A-algebra and f: M → E an A-linear mapping such that

\[
f (x) f (y) = f (y) f (x) \text {   for   all   } x, y \text {   in   } \mathbf {M}.\tag{1}
\]

There exists one and only one A-algebra homomorphism \(g \colon \mathbb{S}(\mathbf{M}) \to \mathbf{E}\) such that \(f = g \circ \phi'\) .

In other words, \((\mathbb{S}(\mathbf{M}), \phi')\) is a solution of the universal mapping problem (Set Theory, IV, §3, no. 1), where \(\Sigma\) is the species of A-algebra structure, the \(\alpha\) -mappings being the linear mappings of the A-module M to an A-algebra satisfying (1).

The uniqueness of \(g\) follows from the fact that \(\phi'(\mathbf{M}) = \mathbf{M}\) generates \(\mathbf{S}(\mathbf{M})\) . To prove the existence of \(g\) , note that by virtue of § 5, no.~1, Proposition 1 there exists an A-algebra homomorphism \(g_1: \mathbf{T}(\mathbf{M}) \to \mathbf{E}\) such that \(f = g_1 \circ \phi\) ;

all that needs to be proved is that \(g_1\) is zero on the ideal \(\mathfrak{J}'\) , for then, if \(p: \mathsf{T}(\mathbf{M}) \to \mathsf{S}(\mathbf{M}) = \mathsf{T}(\mathbf{M}) / \mathfrak{J}'\) is the canonical homomorphism, we can write \(g_1 = g \circ p\) , where \(g: \mathsf{S}(\mathbf{M}) \to \mathbf{E}\) is an algebra homomorphism, and the conclusion will follow from the fact that \(p \circ \phi = \phi'\) . Now the kernel of \(g_1\) is a two-sided ideal which, by virtue of (1) and the relation \(g_1 \circ \phi = f\) , contains the elements \(x \otimes y - y \otimes x\) for \(x, y\) in \(\mathbf{M}\) . This completes the proof.

Remarks. (1) Suppose that E is a graded A-algebra of type Z, with graduation \((\mathrm{E}_{n})\) , and suppose also that the linear mapping f (assumed to satisfy (1)) is such that

\[
f (\mathbf {M}) \subset \mathrm{E} _ {1}.\tag{2}
\]

Then the relation \(g(x_1x_2 \ldots x_p) = f(x_1)f(x_2) \ldots f(x_p)\) with the \(x_i \in \mathbf{M}\) shows that \(g(\mathbf{S}^p(\mathbf{M})) \subset \mathrm{E}_p\) for all \(p \geqslant 0\) and hence \(g\) is a graded algebra homomorphism.

\begin{enumerate}
\def\labelenumi{(\arabic{enumi})}
\setcounter{enumi}{1}
\item
  Every element of \(\mathsf{S}(\mathbf{M})\) is a sum of products of the form \(x_{1}x_{2}\ldots x_{n}\) , where the \(x_{i}\) belong to \(\mathbf{M}\) ; care should be taken not to confuse such products taken in \(\mathsf{S}(\mathbf{M})\) with the analogous products taken in \(\mathsf{T}(\mathbf{M})\) .
\item
  If \(n!\) .1 is invertible in A, the A-module \(S^n(M)\) is generated by the elements of the form \(x^n\) , where \(x \in M\) ; this follows from the above remark and I, § 8, no. 2, Proposition 2.
\end{enumerate}

